\thispagestyle{empty}
\centerline{\textbf{\large{چکیده}}}
\begin{quote}
	در دنیای واقعی، غالبا با مسائل با تعداد دسته زیاد و دادگان کم در هر دسته روبرو هستیم. به بیان دیگر، مسائل در دنیای واقعی دسته‌بندی بزرگ مقیاس با نمونه محدود هستند.
در حالی که پیشرفت و موفقیت مدل‌های یادگیری ژرف عمدتا در مسئله‌هایی با تعداد دسته محدود، ثابت و با دادگان بسیار زیاد برای هر کلاس است. 

روش‌های نمونه محدود در سال‌های اخیر با معرفی رویکردهای نوین در قالب متایادگیری به کارایی قابل توجهی رسیده‌اند. با این حال، هنوز این روش‌ها در تنظیمات مقیاس بزرگ که چند هزار کلاس برای دسته‌بندی وجود دارد، به چالش می‌خورند.

در این پژوهش، مسئله در قالب دو سطح دسته‌بندی با ریزدانگی متفاوت در قالب یک یادگیری سلسله‌مراتبی در یک چارچوب بیزی برای ترکیب اطلاعات مدل شده است. در سطح درشت‌دانه، یک روش تشخیص شی که مقاوم به تغییر زیر جمعیت است، معرفی شده است. 

در سطح ریزدانه، دسته‌بندهایی با رویکرد مولد طراحی شده‌اند که متخصص دسته‌بندی در حیطه‌ی یک کلاس درشت‌دانه هستند. این دسته‌بندها، یادگیرهای نمونه محدودی هستند که در فرآیند یادگیری خود، از اطلاعات سطح درشت‌دانه استفاده می‌کنند و در تنظیمات متایادگیری آموزش می‌بینند تا خواص تعمیم‌پذیری پیدا کنند.

همچنین، از مدل‌های بنیادی پیش‌آموزش دیده برای رسیدن به بازنمایی‌های غنی و بهینگی محاسبات بهره برده شده است. در واقع این پژوهش یک تنظیمات بدیع برای مسئله‌ی یادگیری بزرگ مقیاس و نمونه محدود با وجود چالش‌هایی مانند تغییر زیر جمعیت ارائه داده است. به وسیله‌ی آزمایش‌های مختلف، در تنظیمات با یادگیری و بدون یادگیری نشان داده‌ایم با تولید بازنمایی‌های متاثر از اطلاعات سلسله‌مراتبی می‌توان به ترتیب تا ۱/۳۱ درصد و ۲/۷۵ درصد نسبت به مدل پایه افزایش دقت نسبی داشت.

\vskip 1cm
\textbf{کلمات کلیدی:} \textiranic{
يادگيری ژرف،
یادگیری سلسله‌مراتبی،
دسته‌بندي بزرگ مقياس،
يادگيری بانمونه‌ی محدود،
متايادگيری،
تعمیم‌پذیری
}
\end{quote}